% Section 7 — Limitations and future research
\section{Limitations and Future Research}
\label{sec:limits}

\subsection{Dataset scope}
\label{sec:l-data}
Two retail environments, however contrasted, remain two. Both reflect
store-based daily demand; neither covers e-commerce spikes, perishables with
expiry, promotion-driven regimes, or supply-side censoring beyond the
lost-sales assumption.

\subsection{Series selection}
\label{sec:l-series}
M5 results describe the stratified 500, not the 30,490-series catalog.
Stratification protects archetype coverage but does not guarantee
catalog-level point estimates; extending the full ladder to the catalog is
compute-bound (SARIMA alone required 2,382 seconds for 500 series).

\subsection{Forecast horizon}
\label{sec:l-horizon}
All conclusions are conditional on $h=28$. Shorter horizons tend to favor
naive and averaging methods; longer ones may favor trend and seasonal
specifications. The robustness grid varies policy, not horizon.

\subsection{Inventory-policy scope}
\label{sec:l-policy}
One policy family (daily-review order-up-to, lost sales, no capacity limits,
no minimum order quantities, no perishability, stationary costs) carries all
decision conclusions. Backordering regimes, batch constraints, or stochastic
lead times could reorder the boards. The stability discipline, rather than
any single ranking, is the transferable element.

\subsection{Feature scope}
\label{sec:l-features}
Only univariate demand was used: no prices, promotions, calendar events, or
cross-series covariates. Additional features could widen the LSTM's
advantage or improve the classical methods with regressors. Both directions
remain open questions.

\subsection{Archetype limitations}
\label{sec:smooth-cap}
The Smooth archetype contains exactly one series ($n=1$). Nothing in this
report generalizes from it, and its figures appear for completeness only.

\subsection{SARIMA scope}
\label{sec:l-sarima}
SARIMA is Store-only by design; no claim is made about seasonal
autoregression on intermittent demand. The 100-series subset is exploratory
history, not the headline result.

\subsection{Computational considerations}
\label{sec:l-compute}
Per-estimation costs (approximately 0.03 seconds for ARIMA and 0.6 seconds
for SARIMA, plus LSTM training) bound full-catalog scaling. Per-series ARIMA
order search was excluded on this ground; larger budgets could revisit it.

\subsection{Statistical dependence}
\label{sec:l-dependence}
As stated in \cref{sec:stats-method}, the 112,000 paired errors per
comparison are repeated observations from 500 series across overlapping
origins and horizons, not independent draws. Reported significance levels
therefore describe the preregistered paired design, not inference at a
112,000-strong independent sample size. This limitation applies symmetrically
to all pairwise verdicts in the report.

\subsection{LLM experiment --- future work only}
\label{sec:llm-future}
The program's final rung (local LLM forecasting with controlled versus
context-enhanced prompting, leak-guarded prompt construction in
\texttt{11\_src/llm\_experiment\_design.md}) was designed but deliberately
not executed, held behind a frozen benchmark so the benchmark could not be
adjusted against it. \textbf{No LLM result exists; none is presented,
implied, or charted anywhere in this report.} Executing that phase against
the frozen evidence base is the natural next step.

\subsection{Proposed future extensions}
\label{sec:l-future}
In priority order: (1) execute the LLM phase; (2) add price, promotion, and
calendar covariates; (3) widen the policy family (backorders, minimum order
quantities, perishability); (4) vary the horizon (7 and 56 days); (5) extend
the inexpensive rungs toward the full M5 catalog; (6) test additional retail
environments; (7) replicate the pairwise inference with series-level and
block-bootstrap procedures that model the dependence in \cref{sec:l-dependence}
explicitly.
